\subsection{محدودیت‌ها و تحلیل شکست}
تحلیل مکانیزم عملکردی و خطاهای تجربی ماژول نکات زیر را آشکار می‌سازد:
\begin{itemize}
    \item \textbf{حساسیت بالا به پوشش مجموعه داده کالیبراسیون:} تقطیر وزن‌ها بر پایه ماتریس هم‌کوواریانس تجربی $\mathbf{C}_{YX}$ وابستگی مستقیمی به توزیع آماری داده‌های ورودی کالیبراسیون دارد؛ در صورت محدود بودن دسته‌ها، ماتریس هم‌کوواریانس دچار سوگیری شده و توان تعمیم مدل کاهش می‌یابد.
    \item \textbf{انباشت خطا در فرآیند تولید اتورگرسیو:} در مدل‌های دیکودر لایه‌به‌لایه‌محور (\lr{GPT-2} و \lr{Qwen})، تقریب وزن‌ها با بازنویسی مستقیم وزن بدون فاکتورگیری تفکیک‌شده، خطاهای کوچکی در گام‌های اولیه بازتولید می‌کند که در گذر گام‌های متوالی تولید بازگشتی بزرگ‌نمایی شده و افت دقت را تشدید می‌نماید.
    \item \textbf{مصرف حافظه حین فاز جمع‌آوری فعال‌سازی‌ها:} تجمیع تانسورهای ورودی و خروجی تمام لایه‌ها در رم سیستم نیازمند تخصیص پویای منابع و تخلیه مکرر بافرها در مدل‌های بسیار عمیق است.
\end{itemize}
